% CITED INPUT ONLY night:P1.1
\begin{theorem}\label{night:P1.1}
\StatementRecord{night:P1.1}{proved}{}{Definitions and the displayed constructions}
The finite homology monodromies have orders $5,4$, respectively, and
characteristic polynomials
$(T-1)^2\Phi _5(T)$ and $(T-1)^2(T+1)^2(T^2+1)$.
For $N=M_0-I$ we have $N^2=0$, $\operatorname{rank}N=3$, and
nonzero Smith factors $1,1,1$.  Moreover
\[
 \Lambda^{\langle M_1,M_2\rangle}=0,\qquad
 \Lambda_{\langle M_1,M_2\rangle}=\mathbb Z,
\]
where the quotient homomorphism is $\ell(x)=x_0$.
On the cohomology lattice the corresponding groups are $\mathbb Ze_0$
and $\mathbb Z/2$.  Every nonidentity power of each specified affine
action $\tau_i$ is free.
\end{theorem}
\begin{proof}
\begin{enumerate}[label=\arabic*.]
\item Exact multiplication gives
\[
N=\begin{pmatrix}
0&0&0&0&0&0\\0&0&0&1&0&-1\\0&0&0&0&1&1\\
-1&0&0&0&0&0\\1&0&0&0&0&0\\-1&0&0&0&0&0
\end{pmatrix},\qquad
S=\begin{pmatrix}
0&0&0&1&0&0\\0&1&0&0&0&0\\0&0&1&0&0&0\\
-1&0&0&0&1&0\\1&0&0&0&0&1\\-1&0&0&0&0&0
\end{pmatrix}.
\]
The matrix $S$ is unimodular and
$S^{-1}M_0S=\left(\begin{smallmatrix}I_3&I_3\\0&I_3\end{smallmatrix}\right)$.
The columns of $S$ are henceforth denoted
$(a_1,a_2,a_3,b_1,b_2,b_3)$, so that $Nb_j=a_j$.
This is also an integral certificate of the Smith assertion.
The Hermite basis of
$(M_1-I)\Lambda+(M_2-I)\Lambda$ is $(e_1,\ldots,e_5)$;
the stacked matrix of the two fixed-vector equations has rank six.
Transposing the matrices gives a fixed lattice $\mathbb Ze_0$ and
coinvariant Smith factors $(1,1,1,1,1,2)$.  The pairing
$\Lambda^\vee\times\Lambda\longrightarrow\mathbb Z$ identifies the
fixed cohomology vector $e_0$ with the covector $\ell$.

\item We have $M_iv_i=v_i$, $\ell(v_1)=1$, and $\ell(v_2)=-1$.
Thus $\tau_i^k(x)=M_i^kx+kv_i/m_i$, and $\tau_i^{m_i}=1$ on the
torus.  A fixed point of $\tau_i^k$ would give
$k\ell(v_i)/m_i=0$ in $\mathbb R/\mathbb Z$.  For
$0<k<m_i$ this is impossible.  The orders, characteristic polynomials, fixed-space dimensions,
coinvariant Smith factors, primitive cusp image and every affine-power
obstruction are recomputed exactly in
\path{code/current_audits/audit_corrections/code/independent_V_connection_spectra_06.py}.
Its JSON output retains each assertion. The displayed matrix $S$ and its
standard-form equality are also verified in
\path{code/current_audits/audit_corrections/code/independent_V_markings_06.py}.

\end{enumerate}
\end{proof}

% CITED INPUT ONLY night:P1.9
\begin{theorem}\label{night:P1.9}
\StatementRecord{night:P1.9}{proved}{}{Definitions and the displayed constructions}
The periodic Kuhn fan, in the marking $S$, gives a smooth proper
holomorphic filling of the periods $Z(q)=sI_3+C(q)$ for a sufficiently
small disk.  Its central fibre is reduced with normal crossings.
Modulo the index-one period lattice it has respectively $1,7,12,6$
strata of dimensions $3,2,1,0$, each isomorphic to the corresponding
complex torus $(\mathbb C^*)^d$, and Euler number $6$.
An index-$d$ period sublattice gives $d$ times these stratum counts and
Euler number $6d$.
\end{theorem}
\begin{proof}
\begin{enumerate}[label=\arabic*.]
\item Apply \cite[Theorem~8.7 and Corollary~8.8, attached-source edition, pp.~57--59]{Yang2026}
with $n=3$, $t=q$, $D(q)=C(q)$, $B=I_3$, and the displayed
periodic Kuhn triangulation.  These substitutions give exactly the
action
\[
 (z,q)\longmapsto
 \bigl(q^\lambda\exp(2\pi iC(q)\lambda)z,q\bigr)
\]
and the period order $[I_3\mid sI_3+C(q)]$, where
$q=\exp(2\pi i s)$.  Thus the identity on the dense torus, with
the same integral monomial charts, identifies our quotient with
the cited quotient.  The source theorem supplies its smoothness,
Hausdorffness, properness, and the reduced local equations
$q=x_0\cdots x_k$.  Flatness follows since this nonconstant
function is a nonzero divisor in each smooth chart.
\item For an index-$d$ period sublattice $\Gamma=B\mathbb Z^3$,
the same substitutions give the source's quotient map induced
by $\Gamma\subset\mathbb Z^3$.  It has degree $d$, and
Corollary~8.8 gives $\chi(W)=3!d=6d$ for this marked Kuhn fan.
\end{enumerate}
\end{proof}
